\documentclass{ltc-guji}
\setmainfont{TW-Kai}
\关闭分页
\无标点模式

% 复现 bug: \按 中连续多个 \平抬/\单抬 后，最后的 \平抬 缩进不为0
% 模拟四库全书简明目录的结构

\NewDocumentEnvironment{按环境}{}
  {\begin{段落}[indent=4]\begin{夹注环境}[auto-balance=false]}
  {\end{夹注环境}\end{段落}}

\NewDocumentCommand{\按}{ +m }{%
  \begin{按环境}
    #1%
  \end{按环境}
}

\begin{document}
\begin{正文}
《子夏易傳》十一卷

\按{謹按：唐徐堅《初學記》以太宗御制升列歷代之前，蓋尊尊之大義。焦竑《國史經籍志》，朱彞尊《經義考》並踵前規。臣等編摩《四庫》，初亦恭
\单抬《御定易經通注》
\单抬《御纂周易折中》
\平抬《御纂周易述義》，弁冕諸經。仰蒙
\平抬 指示，命冠於
\相对抬头[3]{國朝} 著述之首，俾尊卑有序，而時代不淆。\平抬 聖度謙沖，酌中立憲，實為千古之大公。謹恪遵\平抬 彞訓，仍託始於《子夏易傳》。並發凡於此，著《四庫》之通例}
\按{焉。}

後續正文
\end{正文}
\end{document}
